\documentclass[12pt,a4paper]{article}

\usepackage[T1]{fontenc}
\usepackage[utf8]{inputenc}
\usepackage{geometry}
\usepackage{hyperref}
\usepackage{longtable}
\usepackage{enumitem}
\usepackage{xcolor}
\usepackage{graphicx}
\usepackage{float}
\usepackage{url}

\geometry{margin=2.5cm}
\hypersetup{
    colorlinks=true,
    linkcolor=blue,
    urlcolor=blue
}

\title{Sudoku Game\\Instrukcja użytkownika}
\author{}
\date{\today}

\begin{document}
\maketitle
\tableofcontents
\newpage

\section{Przeznaczenie programu}

Sudoku Game jest konsolową grą logiczną napisaną w języku C++.
Program pozwala rozpocząć nową rozgrywkę Sudoku, wybrać poziom trudności,
zmienić rozmiar planszy, zapisać stan gry oraz wczytać wcześniejszy zapis.

Aplikacja działa w terminalu. Do sterowania używa klawiatury, dlatego najlepiej
uruchamiać ją w zwykłym oknie terminala, a nie w panelu wyjścia IDE.

\section{Wymagania}

\begin{itemize}[noitemsep]
    \item System Linux albo Windows.
    \item Terminal obsługujący sekwencje ANSI.
    \item Dla samodzielnej kompilacji: CMake oraz kompilator C++17.
\end{itemize}

Gotowy plik wykonywalny dla Windows znajduje się po zbudowaniu pod ścieżką
\path{build/windows-x64-static/Sudoku_Game.exe}. Wariant statyczny jest
najwygodniejszy do przenoszenia na inny komputer z Windows.

\section{Uruchamianie}

\subsection{Linux}

\begin{verbatim}
cmake -S . -B build
cmake --build build
./build/Sudoku_Game
\end{verbatim}

\subsection{Windows}

Jeżeli dostępny jest gotowy plik \texttt{Sudoku\_Game.exe}, należy uruchomić go
w terminalu Windows, PowerShell albo Windows Terminal. Program jest aplikacją
konsolową, więc nie otwiera osobnego okna graficznego.

\section{Menu główne}

Po uruchomieniu programu wyświetlane jest menu. Aktywna opcja jest wyróżniona
kolorem tła.

\begin{longtable}{p{0.3\textwidth}p{0.6\textwidth}}
\textbf{Opcja} & \textbf{Opis} \\
\hline
Nowa gra & Tworzy nową planszę zgodnie z aktualnym trybem i poziomem trudności. \\
Wczytaj zapis & Wczytuje ostatni zapis z pliku \texttt{sudoku\_save.txt}. \\
Poziom trudności & Przełącza poziomy: łatwy, średni i trudny. \\
Tryb gry & Przełącza planszę 9x9 albo 16x16. \\
Wyjście & Kończy działanie programu. \\
\end{longtable}

\section{Sterowanie}

\begin{longtable}{p{0.3\textwidth}p{0.6\textwidth}}
\textbf{Klawisz} & \textbf{Działanie} \\
\hline
Strzałki & Przesuwają zaznaczone pole albo wybór w menu. \\
Enter & Zatwierdza opcję w menu. W trybie 16x16 zatwierdza wartość 1, jeżeli wcześniej naciśnięto klawisz 1. \\
1--9 & Wpisują wartość do pola gracza. \\
0 & Czyści wybrane pole gracza. \\
Backspace, Delete & Czyści wybrane pole gracza. \\
s & Zapisuje aktualny stan gry. \\
q & Wraca z gry do menu albo zamyka program z poziomu menu. \\
\end{longtable}

\section{Zasady gry}

Celem gry jest uzupełnienie całej planszy tak, aby w każdym wierszu, każdej
kolumnie i każdym bloku występowała każda dopuszczalna wartość dokładnie raz.
Program nie pozwala wpisać liczby, która łamie zasady Sudoku.

Pola startowe są stałe i nie można ich zmienić ani wyczyścić. Pola wpisane przez
gracza można poprawiać w dowolnym momencie.

\section{Tryby gry}

\subsection{Sudoku 9x9}

Klasyczny wariant Sudoku. Plansza ma 9 wierszy, 9 kolumn i bloki 3x3. Do pól
można wpisywać wartości od 1 do 9.

\subsection{Sudoku 16x16}

Większy wariant gry. Plansza ma 16 wierszy, 16 kolumn i bloki 4x4. Do pól można
wpisywać wartości od 1 do 16.

Wartości od 10 do 16 wpisuje się jako sekwencję dwóch klawiszy:
\begin{itemize}[noitemsep]
    \item \texttt{1}, potem \texttt{0} wpisuje 10,
    \item \texttt{1}, potem \texttt{1} wpisuje 11,
    \item \texttt{1}, potem \texttt{6} wpisuje 16.
\end{itemize}

Jeżeli użytkownik chce wpisać samo 1 w trybie 16x16, powinien nacisnąć
\texttt{1}, a następnie \texttt{Enter}.

\section{Zapis i odczyt gry}

Klawisz \texttt{s} zapisuje aktualną rozgrywkę do pliku
\texttt{sudoku\_save.txt} w katalogu, z którego uruchomiono program. Zapis
przechowuje:

\begin{itemize}[noitemsep]
    \item rozmiar planszy,
    \item poziom trudności,
    \item wartości pól,
    \item informację o polach startowych.
\end{itemize}

Opcja \textit{Wczytaj zapis} odtwarza ten plik. Jeżeli pliku nie ma albo ma
niepoprawny format, program wraca do menu i wyświetla komunikat o błędzie.

\section{Kolory w programie}

\begin{longtable}{p{0.3\textwidth}p{0.6\textwidth}}
\textbf{Kolor} & \textbf{Znaczenie} \\
\hline
Biały & Pole startowe, którego nie można zmienić. \\
Zielony & Wartość wpisana przez gracza. \\
Szary & Puste pole. \\
Niebieskie tło & Aktualnie zaznaczone pole. \\
\end{longtable}

\section{Zakończenie rozgrywki}

Po poprawnym uzupełnieniu całej planszy program wyświetla komunikat
gratulacyjny. Naciśnięcie dowolnego klawisza powoduje powrót do menu.

\end{document}
